\begin{proofbox}
	\textbf{\textcolor{CProof}{思路提示}}
	平行转化链的四个定理可以采用反证法或直接利用平行定义证明。核心思路是：通过假设反例导出矛盾，或利用平行线的唯一性。

	\medskip
	\textbf{\textcolor{CProof}{正式推导}}
	\begin{description}[style=nextline, leftmargin=2.8em, labelsep=0.5em, itemsep=0.55em]
		\item[\textbf{步骤一（线面平行判定）：}]
			\[
				a\parallel b,\ a\not\subset\alpha,\ b\subset\alpha \Rightarrow a\parallel\alpha
			\]
			\par\textbf{理由：} 反证法。假设$a$与$\alpha$相交于点$A$，则$A\in\alpha$。在平面$\alpha$内过$A$作直线$c\parallel b$。因为$a\parallel b$，且$a$与$c$都过$A$且平行于$b$，根据平行线的唯一性，$a$与$c$重合，从而$a\subset\alpha$，与$a\not\subset\alpha$矛盾。故$a\parallel\alpha$。
		\item[\textbf{步骤二（面面平行判定）：}]
			\[
				a,b\subset\alpha,\ a\cap b=P,\ a\parallel\beta,\ b\parallel\beta \Rightarrow \alpha\parallel\beta
			\]
			\par\textbf{理由：} 反证法。假设$\alpha\cap\beta=l$。因为$a\parallel\beta$，且$a\subset\alpha$，所以$a$与$l$在同一平面$\alpha$内。若$a$与$l$相交，则交点在$l$上也在$a$上，从而在$\beta$上，与$a\parallel\beta$矛盾。所以$a\parallel l$。同理$b\parallel l$。于是$a\parallel b$，与$a\cap b=P$矛盾。故$\alpha\parallel\beta$。
		\item[\textbf{步骤三（面面平行性质）：}]
			\[
				\alpha\parallel\beta,\ \gamma\cap\alpha=a,\ \gamma\cap\beta=b \Rightarrow a\parallel b
			\]
			\par\textbf{理由：} 直接法。因为$\alpha\parallel\beta$，所以$\alpha$与$\beta$没有公共点。直线$a$和$b$都在平面$\gamma$内。若$a$与$b$不平行，则它们相交于一点$Q$。点$Q$既在$a$上又在$b$上，所以$Q\in\alpha$且$Q\in\beta$，与$\alpha\parallel\beta$矛盾。因此$a\parallel b$。
		\item[\textbf{步骤四（线面平行性质）：}]
			\[
				a\parallel\alpha,\ a\subset\beta,\ \alpha\cap\beta=b \Rightarrow a\parallel b
			\]
			\par\textbf{理由：} 反证法。假设$a$与$b$不平行，由于它们同在平面$\beta$内，故相交于一点$Q$。$Q\in a$且$Q\in b$。因为$b\subset\alpha$，所以$Q\in\alpha$；又$Q\in a$，所以$a$与$\alpha$有公共点$Q$，与$a\parallel\alpha$矛盾。故$a\parallel b$。
	\end{description}

	\medskip
	\textbf{\textcolor{CProof}{结论回扣}}
	上述四个定理构成了平行关系的完整转化链：步骤一是从线线平行到线面平行，步骤二从线面平行到面面平行，步骤三、四则从面面平行、线面平行推出线线平行。它们相辅相成，是立体几何平行证明的基础。
\end{proofbox}
